% paper/sections/09b2_split_ppe.tex
%
% §9.3  分相 PPE 解法：密度比非依存な高精度 PPE 戦略
% ═══════════════════════════════════════════════════════════════════════════════
\subsection{分相 PPE 解法：密度比非依存な高精度 PPE 戦略}
\label{sec:split_ppe_framework}
% backward-compat labels (previously in 12e_interface_crossing.tex)
\label{sec:split_ppe_motivation}%
\label{sec:gfm_numerical_motivation}%
% ═══════════════════════════════════════════════════════════════════════════════

\noindent\textbf{本節の位置づけ：}
§~\ref{sec:ccd_poisson_matrix} で構築した CCD 離散化は
等密度ポアソン方程式に対して $\Ord{h^6}$ の設計精度を達成する．
しかし二相流の PPE は変密度演算子 $\bnabla\cdot(\rho^{-1}\bnabla p) = q$ であり，
界面を横切るステンシルが $\Ord{1}$ の打ち切り誤差を生じる．
本節では，この本質的制約を回避する\textbf{分相 PPE 解法}を本稿の主要 PPE 戦略として定式化する．

% ─────────────────────────────────────────────────────────────────────────────
\subsubsection{変密度一括解法の本質的制約}
\label{sec:monolithic_ppe_limitation}
% ─────────────────────────────────────────────────────────────────────────────

Smoothed Heaviside による一括解法では，遷移層（幅 $\varepsilon = 1.5h$）内で
$1/\rho$ が $\rho_l/\rho_g$ 倍のジャンプを持つ．
界面を横切る CCD/FD ステンシルは $\Ord{1}$ の打ち切り誤差を生じ，
PPE のグリーン関数の非局所性により解全体に伝播する
（§~\ref{sec:ppe_condition_number}）．
この誤差は反復回数 $k$ の増大や緩和係数 $\omega$ では回復不能であり，
smoothed Heaviside 一括解法の\textbf{本質的制約}である：
%
\begin{itemize}[noitemsep]
  \item $\rho_l/\rho_g = 2$--$5$：格子収束次数が $\Ord{h^{0.1}}$--$\Ord{h^{2}}$ に劣化
  \item $\rho_l/\rho_g \geq 10$：DC が発散（反復行列のスペクトル半径 $> 1$）
  \item 水--空気系（$\rho_l/\rho_g \approx 830$）：使用不可
\end{itemize}
%
§~\ref{sec:interface_crossing} に定量的な破綻データを示す．
重要な点として，$\rho_l/\rho_g = 2$--$5$ でも界面近傍の精度は
CCD の設計精度 $\Ord{h^6}$ から\textbf{大幅に劣化}しており，
低密度比だから一括解法で十分ということはない．

% ─────────────────────────────────────────────────────────────────────────────
\subsubsection{分相 PPE 解法の定式化}
\label{sec:split_ppe_formulation}
% ─────────────────────────────────────────────────────────────────────────────

分相 PPE 解法の核心は「各相で定数密度のラプラシアンを解く」ことにある．
Smoothed Heaviside 一括解法が $\bnabla\cdot(\rho^{-1}\bnabla p) = q$ を
変動係数 $\rho(\psi)$ のまま全領域で解くのに対し，
分相 PPE 解法は各相 $k\in\{\ell,g\}$ で
%
\begin{equation}
  \nabla^2 p_k = \rho_k \, q \qquad \text{（$\rho_k$ は定数）}
  \label{eq:split_ppe}
\end{equation}
%
を解く．ただし，これは単に各相を切り離して解くという意味ではない．
界面 $\Gamma$ 上では圧力ジャンプと法線フラックス条件を同時に課して，
二つの相の楕円型問題を閉じる：
\begin{align}
  [p]_\Gamma
    &= J_p
     = \sigma\kappa
     \quad \text{（圧力形式）},
     \label{eq:split_ppe_pressure_jump} \\
  \left[\frac{1}{\rho}\pd{p}{n}\right]_\Gamma
    &= J_n .
     \label{eq:split_ppe_flux_jump}
\end{align}
IPC の圧力増分 $\delta p$ を未知量にする場合は，
$J_p=\sigma(\kappa^{n+1}-\kappa^n)$ とし，
法線速度の連続性から
$J_n=\Delta t^{-1}[\bu^*\cdot\hat{\bm n}]$ を用いる．
予測速度が界面法線方向に連続となる離散化では $J_n=0$ に退化する．
HFE（§~\ref{sec:field_extension}）は，これらのジャンプ条件で補正した
片側データを界面越しに延長し，CCD ステンシルが不連続値に触れないようにするための
補助技術である．
$\rho_k$ が定数であるため，各相の離散化演算子は密度ジャンプの影響を一切受けない．

\noindent\textbf{欠陥補正法との関係：}
式~\eqref{eq:split_ppe} は等密度ポアソン方程式であるから，
§~\ref{sec:defect_correction_main} の欠陥補正法（DC $k = 3$）を各相内で適用できる．
等密度条件の下では DC の CCD 高次演算子 $L_H$ と FD 低次演算子 $L_L$ の
スペクトル不整合が最小化されるため，
§~\ref{sec:dc_accuracy} の理論が示す $\Ord{h^7}$（Dirichlet 超収束）
ないし $\Ord{h^5}$（Neumann BC 律速）が\textbf{密度比によらず}回復する．

\begin{tcolorbox}[resultbox, title={分相 PPE + DC $k{=}3$ の精度保証}]
\begin{enumerate}[noitemsep]
  \item 各相で定数密度の PPE を独立に解くため，
        界面条件~\eqref{eq:split_ppe_pressure_jump}--\eqref{eq:split_ppe_flux_jump}を満たす
        片側問題として閉じる
  \item DC $k{=}3$ は各相内で $\Ord{h^7}$（Dirichlet）/
        $\Ord{h^5}$（Neumann BC 律速）を達成する
  \item 各相内の離散誤差は密度比に直接依存しない．
        ただし全体精度は界面ジャンプ条件と HFE 延長の精度に律速される
        （§~\ref{sec:verify_split_ppe_dc} で検証）
  \item 速度補正の勾配 $\bnabla(\delta p)$ には CCD $\Ord{h^6}$ を適用し，
        Balanced-Force 整合性を保持する
\end{enumerate}
\end{tcolorbox}

\noindent\textbf{HFE の役割：}
分相 PPE 解法において HFE（§~\ref{sec:field_extension}）は不可欠な基盤技術である．
IPC の $\bnabla p^n$（Predictor ステップ）および
$\bnabla(\delta p)$（Corrector ステップ）が界面を横切る前に，
圧力場をジャンプ補正済み HFE Hermite 補間で滑らかに延長することで，
CCD ステンシルが不連続値に触れることを防止する．
